\lecture{06}{Строгая и сильная выпуклость}

\sourcevideo
    {https://www.youtube.com/playlist?list=PLOzRYVm0a65fhKDS8cYIC7YCuVGJ4FOJx}
    {Week 2: Lecture 6: Strictly and Strongly Convex Functions}

Выпуклость гарантирует глобальность локальных минимумов, но не
единственность решения и не задаёт количественной оценки кривизны.
Строгая выпуклость обеспечивает единственность минимума, если он
существует, а сильная выпуклость вводит равномерную квадратичную
нижнюю оценку. Она позволяет получать явные оценки скорости
сходимости.

\section{Строгая выпуклость и условие второго порядка}

Пусть \(C\subseteq\RR^n\) --- выпуклое множество.

\begin{definition}[Строго выпуклая функция]
Функция \(f\colon C\to\RR\) называется строго выпуклой, если для
любых различных \(x,y\in C\) и каждого \(\theta\in(0,1)\)
\[
    f\bigl(\theta x+(1-\theta)y\bigr)
    <
    \theta f(x)+(1-\theta)f(y).
\]
\end{definition}

Для обычной выпуклости допускается равенство. Поэтому выпуклая
функция может иметь плоский участок и несколько точек минимума.
Строгая выпуклость исключает такой участок между различными точками.

\begin{theorem}[Единственность минимума]
Если строго выпуклая функция достигает глобального минимума, то её
точка минимума единственна.
\end{theorem}

\begin{proof}
Предположим, что различные точки \(x^\star\) и \(y^\star\) являются
минимумами:
\[
    f(x^\star)=f(y^\star)=f^\star.
\]
Для \(\theta\in(0,1)\) строгая выпуклость даёт
\[
\begin{aligned}
    f\bigl(\theta x^\star+(1-\theta)y^\star\bigr)
    &<
    \theta f(x^\star)+(1-\theta)f(y^\star)
    \\
    &=f^\star,
\end{aligned}
\]
что противоречит минимальности \(f^\star\).
\end{proof}

Условие существования существенно: строгая выпуклость сама по себе
не гарантирует, что минимум достигается. Например,
\[
    f(x)=\mathrm{e}^x
\]
строго выпукла на \(\RR\), но не имеет точки глобального минимума.



\begin{theorem}[Достаточное условие строгой выпуклости]
Пусть \(C\subseteq\RR^n\) открыто и выпукло, а
\(f\in C^2(C)\). Если
\[
    \nabla^2 f(x)\succ0
    \qquad
    \text{для всех }x\in C,
\]
то \(f\) строго выпукла на \(C\).
\end{theorem}

Положительная определённость означает
\[
    v^{\mathsf T}\nabla^2f(x)v>0
    \qquad
    \text{для всех }v\ne0.
\]

Это условие достаточно, но не необходимо. Рассмотрим
\[
    f(x)=x^4.
\]
Функция строго выпукла и имеет единственный минимум в нуле, однако
\[
    f''(x)=12x^2,
    \qquad
    f''(0)=0.
\]
Следовательно, матрица Гессе не обязана быть положительно определена
в каждой точке строго выпуклой функции.

Для функции
\[
    f(x)=x^2
\]
имеем
\[
    f''(x)=2>0,
\]
поэтому её строгая выпуклость непосредственно следует из критерия
второго порядка.

Таким образом,
\[
    \nabla^2f(x)\succ0
    \quad\Longrightarrow\quad
    f\text{ строго выпукла},
\]
но обратная импликация в общем случае неверна.

\section{Сильная выпуклость и её характеризации}

Строгая выпуклость утверждает, что значение на внутренней точке
отрезка строго меньше значения соответствующей хорды, но не оценивает
величину разности.

\begin{definition}[Сильная выпуклость]
Функция \(f\colon C\to\RR\) называется \(\mu\)-сильно выпуклой,
где \(\mu>0\), если для любых \(x,y\in C\) и
\(\theta\in[0,1]\)
\[
\begin{aligned}
    f\bigl(\theta x+(1-\theta)y\bigr)
    \le{}&
    \theta f(x)+(1-\theta)f(y)
    \\
    &-
    \frac{\mu}{2}
    \theta(1-\theta)\norm{x-y}^2.
\end{aligned}
\]
Число \(\mu\) называется модулем сильной выпуклости.
\end{definition}

Дополнительное квадратичное слагаемое задаёт равномерную нижнюю
оценку кривизны. Чем больше допустимое значение \(\mu\), тем сильнее
функция отклоняется вниз от своих хорд.

Справедлива цепочка импликаций
\[
    \text{сильная выпуклость}
    \Longrightarrow
    \text{строгая выпуклость}
    \Longrightarrow
    \text{выпуклость}.
\]
Обратные импликации неверны.

Функция \(x^2\) сильно выпукла, тогда как \(x^4\) строго выпукла,
но не сильно выпукла на всей прямой. Действительно,
\[
    \inf_{x\in\RR}12x^2=0,
\]
поэтому её кривизну нельзя равномерно отделить от нуля.



\begin{theorem}[Критерий первого порядка]
Пусть \(C\subseteq\RR^n\) открыто и выпукло, а
\(f\in C^1(C)\). Функция \(f\) является \(\mu\)-сильно выпуклой
тогда и только тогда, когда
\[
    f(y)
    \ge
    f(x)
    +
    \inner{\nabla f(x)}{y-x}
    +
    \frac{\mu}{2}\norm{y-x}^2
\]
для всех \(x,y\in C\).
\end{theorem}

Для обычной выпуклой функции имеется только опорная гиперплоскость:
\[
    f(y)
    \ge
    f(x)+\inner{\nabla f(x)}{y-x}.
\]
При сильной выпуклости функция лежит выше этой гиперплоскости как
минимум на величину
\[
    \frac{\mu}{2}\norm{y-x}^2.
\]

\begin{corollary}[Квадратичный рост около минимума]
Пусть \(C\subseteq\RR^n\) открыто и выпукло, а
\(f\in C^1(C)\) является \(\mu\)-сильно выпуклой. Если
\(x^\star\in C\) --- точка минимума, то для всех \(x\in C\)
\[
    f(x)-f^\star
    \ge
    \frac{\mu}{2}\norm{x-x^\star}^2.
\]
\end{corollary}

\begin{proof}
Поскольку \(x^\star\) --- внутренняя точка минимума,
\(\nabla f(x^\star)=0\). Подстановка \(x=x^\star\) в критерий
первого порядка даёт требуемое неравенство.
\end{proof}

Это неравенство связывает ошибку по значению функции с расстоянием
до оптимального решения.



\begin{theorem}[Критерий второго порядка]
Пусть \(C\subseteq\RR^n\) открыто и выпукло, а
\(f\in C^2(C)\). Тогда \(f\) является \(\mu\)-сильно выпуклой на
\(C\) тогда и только тогда, когда
\[
    \nabla^2f(x)\succeq\mu I_n
    \qquad
    \text{для всех }x\in C.
\]
\end{theorem}

Эквивалентно,
\[
    v^{\mathsf T}\nabla^2f(x)v
    \ge
    \mu\norm{v}^2
\]
для всех \(x\in C\) и \(v\in\RR^n\).

В одномерном случае условие принимает вид
\[
    f''(x)\ge\mu.
\]

Для функции
\[
    f(x)=x^2
\]
имеем
\[
    f''(x)=2,
\]
поэтому она \(2\)-сильно выпукла. Она также является
\(\mu\)-сильно выпуклой для любого \(0<\mu\le2\), но максимальный
модуль равен \(2\).

Для квадратичной функции
\[
    f(x)=\frac12x^{\mathsf T}Qx+b^{\mathsf T}x+c,
    \qquad Q=Q^{\mathsf T},
\]
матрица Гессе равна
\[
    \nabla^2f(x)=Q.
\]
Следовательно,
\[
    f\text{ выпукла}
    \Longleftrightarrow
    Q\succeq0,
\]
а
\[
    f\text{ \(\mu\)-сильно выпукла}
    \Longleftrightarrow
    Q\succeq\mu I_n.
\]
Если \(Q\succ0\), наибольший допустимый модуль сильной выпуклости
равен
\[
    \mu_{\max}=\lambda_{\min}(Q)>0.
\]
При \(\lambda_{\min}(Q)=0\) функция выпукла, но не сильно выпукла на
всём пространстве.



Из критерия первого порядка следует важное свойство градиента.

\begin{theorem}
Если \(f\in C^1(C)\) является \(\mu\)-сильно выпуклой, то
\[
    \inner{\nabla f(x)-\nabla f(y)}{x-y}
    \ge
    \mu\norm{x-y}^2
\]
для всех \(x,y\in C\).
\end{theorem}

\begin{proof}
Сильная выпуклость даёт
\[
    f(y)
    \ge
    f(x)
    +
    \inner{\nabla f(x)}{y-x}
    +
    \frac{\mu}{2}\norm{y-x}^2
\]
и после перестановки \(x,y\)
\[
    f(x)
    \ge
    f(y)
    +
    \inner{\nabla f(y)}{x-y}
    +
    \frac{\mu}{2}\norm{x-y}^2.
\]
Складывая неравенства, получаем
\[
    \inner{\nabla f(x)-\nabla f(y)}{x-y}
    \ge
    \mu\norm{x-y}^2.
\]
\end{proof}

При \(y=x^\star\) и \(\nabla f(x^\star)=0\)
\[
    \inner{\nabla f(x)}{x-x^\star}
    \ge
    \mu\norm{x-x^\star}^2.
\]
По неравенству Коши--Буняковского отсюда следует
\[
    \norm{\nabla f(x)}
    \ge
    \mu\norm{x-x^\star}.
\]

Следовательно, градиент сильно выпуклой функции количественно
контролирует расстояние до минимума. Строгая выпуклость сама по себе
такой оценки не даёт: для \(f(x)=x^4\) градиент \(4x^3\) быстро
исчезает около нуля.

\section{Негладкие выпуклые функции}

Для негладкой функции градиент может не существовать. В этом случае
используется понятие субградиента.

\begin{definition}[Субградиент]
Пусть \(f\colon C\to\RR\) выпукла. Вектор \(g\in\RR^n\) называется
субградиентом \(f\) в точке \(x\in C\), если
\[
    f(y)
    \ge
    f(x)+\inner{g}{y-x}
\]
для всех \(y\in C\).
Множество всех субградиентов обозначается через
\[
    \partial f(x).
\]
\end{definition}

Если \(f\) дифференцируема в \(x\), то
\[
    \partial f(x)=\set{\nabla f(x)}.
\]

Для функции
\[
    f(x)=\abs{x}
\]
имеем
\[
    \partial f(x)
    =
    \begin{cases}
        \set{1}, & x>0,\\
        [-1,1], & x=0,\\
        \set{-1}, & x<0.
    \end{cases}
\]

В точке \(x=0\) обычная производная не существует, но каждая прямая
с наклоном \(g\in[-1,1]\) лежит не выше графика функции.

\begin{definition}[Сильная выпуклость через субградиенты]
Выпуклая функция \(f\) является \(\mu\)-сильно выпуклой, если для
любых \(x,y\in C\) и каждого \(g\in\partial f(x)\)
\[
    f(y)
    \ge
    f(x)+\inner{g}{y-x}
    +
    \frac{\mu}{2}\norm{y-x}^2.
\]
\end{definition}

Это определение позволяет анализировать сильно выпуклые функции,
содержащие негладкие слагаемые.

\section{Неравенство Поляка--Лоясиевича}

\begin{definition}[Неравенство Поляка--Лоясиевича]
Дифференцируемая функция \(f\colon\RR^n\to\RR\) удовлетворяет
неравенству Поляка--Лоясиевича с константой \(\mu>0\), если
\[
    \frac12\norm{\nabla f(x)}^2
    \ge
    \mu\bigl(f(x)-f^\star\bigr)
\]
для всех \(x\), где
\[
    f^\star=\inf_x f(x).
\]
\end{definition}

Это условие также называют условием градиентного доминирования.
Оно утверждает, что большой разрыв по значению функции сопровождается
достаточно большим градиентом.

\begin{theorem}
Всякая дифференцируемая \(\mu\)-сильно выпуклая функция
удовлетворяет неравенству Поляка--Лоясиевича с той же константой
\(\mu\).
\end{theorem}

\begin{proof}
Из сильной выпуклости для любого \(y\)
\[
    f(y)
    \ge
    f(x)
    +
    \inner{\nabla f(x)}{y-x}
    +
    \frac{\mu}{2}\norm{y-x}^2.
\]
Беря инфимум по \(y\) в обеих частях, получаем
\[
    f^\star
    \ge
    \inf_y
    \left\{
        f(x)
        +
        \inner{\nabla f(x)}{y-x}
        +
        \frac{\mu}{2}\norm{y-x}^2
    \right\}.
\]
Минимум выражения в фигурных скобках достигается при
\[
    y=x-\frac1\mu\nabla f(x)
\]
и равен
\[
    f(x)-\frac{1}{2\mu}\norm{\nabla f(x)}^2.
\]
После перестановки слагаемых получаем PL-неравенство.
\end{proof}

Обратное утверждение неверно: функция может удовлетворять
неравенству Поляка--Лоясиевича, не будучи сильно выпуклой и даже
не будучи выпуклой.

В лекции в качестве такого примера приводится функция
\[
    f(x)=x^2+3\sin^2x.
\]
Она имеет единственный глобальный минимум в нуле и не является
выпуклой на всей прямой, поскольку
\[
    f''(x)=2+6\cos(2x)
\]
принимает отрицательные значения. Для включения функции в класс PL
требуется отдельная глобальная оценка
\[
    \frac12\lvert f'(x)\rvert^2
    \ge
    \mu f(x).
\]
В лекции такая оценка не доказывается, поэтому пример фиксируется как
заявленный, но не доказанный.

Корректное соотношение классов имеет вид
\[
    \{\text{сильно выпуклые функции}\}
    \subset
    \{\text{PL-функции}\},
\]
а также
\[
    \{\text{сильно выпуклые}\}
    \subset
    \{\text{строго выпуклые}\}
    \subset
    \{\text{выпуклые}\}.
\]
При этом класс PL-функций пересекается с классом выпуклых функций,
но не содержится в нём и не совпадает с ним.

\begin{remark}
Класс PL-функций не образует одну цепочку вложений с классами
выпуклых, строго выпуклых и сильно выпуклых функций. Существуют как
выпуклые, но не сильно выпуклые PL-функции, так и невыпуклые
PL-функции.
\end{remark}

\section{Аффинная композиция}

Пусть \(f\colon\RR^m\to\RR\) является \(\mu\)-сильно выпуклой и
\[
    g(x)=f(Ax+b),
    \qquad
    A\in\RR^{m\times n}.
\]
Для любых \(x,y\in\RR^n\) и \(\theta\in[0,1]\)
\[
\begin{aligned}
    g\bigl(\theta x+(1-\theta)y\bigr)
    \le{}&
    \theta g(x)+(1-\theta)g(y)
    \\
    &-
    \frac{\mu}{2}\theta(1-\theta)
    \norm{A(x-y)}^2.
\end{aligned}
\]
Если \(A\) имеет полный столбцовый ранг, то
\[
    \norm{Av}
    \ge
    \sigma_{\min}(A)\norm{v},
\]
и потому \(g\) является
\(\mu\sigma_{\min}^2(A)\)-сильно выпуклой.

Если \(\ker A\ne\{0\}\), то для некоторого \(v\ne0\)
\[
    g(x+tv)=g(x)
    \qquad
    \text{для всех }t\in\RR.
\]
Следовательно, композиция на всём пространстве не может быть строго,
а значит, и сильно выпуклой.

\begin{remark}[Уточнение рангового условия]
Для композиции \(x\mapsto f(Ax+b)\), определённой на \(\RR^n\),
необходима инъективность отображения \(x\mapsto Ax\), то есть полный
столбцовый ранг \(A\) или, эквивалентно, \(\ker A=\{0\}\).
\end{remark}

Вырожденная квадратичная функция всё же может удовлетворять
PL-неравенству. Пусть
\[
    g(x)=\frac12\norm{Ax-b}^2
\]
и \(x^\star\) --- любое решение задачи наименьших квадратов. Тогда
\[
    g(x)-g^\star
    =
    \frac12\norm{A(x-x^\star)}^2,
\]
а
\[
    \nabla g(x)=A^{\mathsf T}A(x-x^\star).
\]
Если \(A\ne0\), то
\[
    \frac12\norm{\nabla g(x)}^2
    \ge
    \sigma_{\min,+}^2(A)\bigl(g(x)-g^\star\bigr),
\]
где \(\sigma_{\min,+}(A)\) --- наименьшее положительное сингулярное
число \(A\). При нетривиальном ядре минимум не единственен, но
PL-неравенство сохраняется.

\section{Практическое значение}

Строгая выпуклость полезна прежде всего для доказательства
единственности решения. Она не задаёт равномерной скорости роста
функции около минимума.

Сильная выпуклость даёт количественные оценки:
\[
    f(x)-f^\star
    \ge
    \frac{\mu}{2}\norm{x-x^\star}^2,
\]
\[
    \norm{\nabla f(x)}
    \ge
    \mu\norm{x-x^\star},
\]
\[
    \frac12\norm{\nabla f(x)}^2
    \ge
    \mu\bigl(f(x)-f^\star\bigr).
\]
Эти соотношения используются при доказательстве линейной и
экспоненциальной сходимости оптимизационных методов.

PL-неравенство сохраняет последнюю оценку без требования
выпуклости. Поэтому доказательства сходимости, использующие только
PL-условие и общие предположения о гладкости, применимы ко всем
сильно выпуклым задачам и одновременно охватывают некоторые
невыпуклые модели.

Единственность глобального минимума сама по себе не характеризует инвексность, поэтому этот класс функций далее не используется.


Строгая выпуклость гарантирует единственность минимума, если он
существует. Положительная определённость матрицы Гессе достаточна, но
не необходима для строгой выпуклости. Сильная выпуклость задаёт
равномерную квадратичную кривизну:
\[
    \nabla^2f(x)\succeq\mu I_n.
\]
Она влечёт строгую выпуклость, квадратичный рост, сильную
монотонность градиента и PL-неравенство. PL-условие может выполняться
и для невыпуклых функций. Для негладких функций аналогом градиента
служит субградиент.